
\documentclass{article}
\usepackage{colortbl}
\usepackage{makecell}
\usepackage{multirow}
\usepackage{supertabular}

\begin{document}

\newcounter{utterance}

\centering \large Interaction Transcript for game `referencegame', experiment `line\_grids\_rows', episode 2 with EuroLLM{-}22B{-}Instruct{-}2512{-}t1.0.
\vspace{24pt}

{ \footnotesize  \setcounter{utterance}{1}
\setlength{\tabcolsep}{0pt}
\begin{supertabular}{c@{$\;$}|p{.15\linewidth}@{}p{.15\linewidth}p{.15\linewidth}p{.15\linewidth}p{.15\linewidth}p{.15\linewidth}}
    \# & $\;$A & \multicolumn{4}{c}{Game Master} & $\;\:$B\\
    \hline

    \theutterance \stepcounter{utterance}  
    & & \multicolumn{4}{p{0.6\linewidth}}{
        \cellcolor[rgb]{0.9,0.9,0.9}{
            \makecell[{{p{\linewidth}}}]{
                \texttt{\tiny{[P1$\langle$GM]}}
                \texttt{You are given three grids, where each of them is 5 by 5 in size.} \\
\texttt{Grids have empty cells marked with "▢" and filled cells marked with "X".} \\
\texttt{Your task is to generate a referring expression that best describes the target grid while distinguishing it from the two other distractor grids.} \\
\texttt{The first grid is the target grid, and the following two grids are the distractors.} \\
\\ 
\texttt{Target grid:} \\
\\ 
\texttt{X X X X X} \\
\texttt{□ □ □ □ □} \\
\texttt{□ □ □ □ □} \\
\texttt{X X X X X} \\
\texttt{□ □ □ □ □} \\
\\ 
\texttt{Distractor grid 1:} \\
\\ 
\texttt{□ □ □ X X} \\
\texttt{□ □ □ X X} \\
\texttt{□ □ □ X X} \\
\texttt{□ □ □ X X} \\
\texttt{□ □ □ X X} \\
\\ 
\texttt{Distractor grid 2:} \\
\\ 
\texttt{□ □ X X X} \\
\texttt{□ □ X X X} \\
\texttt{□ □ X X X} \\
\texttt{□ □ X X X} \\
\texttt{□ □ X X X} \\
\\ 
\texttt{Instruction: Describe the target grid.} \\
\texttt{Generate the referring expression starting with the tag "Expression: " for the given target grid. Omit any other text.} \\
            }
        }
    }
    & \\ \\

    \theutterance \stepcounter{utterance}  
    & \multicolumn{4}{p{0.6\linewidth}}{
        \cellcolor[rgb]{0.9,0.9,0.9}{
            \makecell[{{p{\linewidth}}}]{
                \texttt{\tiny{[P1$\rangle$GM]}}
                \texttt{The target grid has a solid row of "X" cells in the top row, with the middle row completely filled with "□" (empty cells). It also features a solid row of "X" cells in the bottom row, and the leftmost and rightmost columns are filled with "□" (except for the corner cells, which are "X" but not part of the specified pattern). Additionally, the third row contains two "X" cells located at columns 2 and 4. The second row is entirely empty ("□"). } \\
\\ 
\texttt{The two distractor grids differ significantly: Grid 1 lacks the symmetric row patterning and has distinct corner cells; Grid 2 has row{-}wise solid "X" patterns from top to bottom. } \\
\\ 
\texttt{The distinguishing features collectively are: } \\
\texttt{1. Three{-}layered row pattern (first row: top{-}bound "X", middle: center "□", bottom: bottom{-}bound "X")  } \\
\texttt{2. Per{-}column "□" in middle row except corners  } \\
\texttt{3. Third row with "X" at columns 2 and 4  } \\
\texttt{4. Second row entirely empty  } \\
\texttt{5. Top row corners have "X"  } \\
\texttt{6. Bottom row corners are "X"  } \\
\\ 
\texttt{Thus, the unique grid structure is described by the referring expression:  } \\
\\ 
\texttt{**Expression:** Grid with bottom{-}row "X" and top{-}row "X", middle row per{-}column "□" (corners excluded), third row "X" at columns 2 and 4, and second row entirely "□".} \\
            }
        }
    }
    & & \\ \\

    \theutterance \stepcounter{utterance}  
    & & & \multicolumn{2}{p{0.3\linewidth}}{
        \cellcolor[rgb]{0.9,0.9,0.9}{
            \makecell[{{p{\linewidth}}}]{
                \texttt{\tiny{[GM$|$GM]}}
                \texttt{Invalid generated expression} \\
            }
        }
    }
    & & \\ \\

\end{supertabular}
}

\end{document}
